\section*{要旨}
\addcontentsline{toc}{section}{要旨}

本研究では、イベントやアトラクション運用における混雑緩和および待ち時間試算の自動化を目的とし、監視カメラ等の映像データから待機列に存在する人物およびそのグループ単位をリアルタイムかつ高精度に検出する手法について検討する。待機列内では、人物同士の密集による遮蔽、列の前後に跨るコミュニケーション、一時的な通行人の横切りなど、単一のオブジェクト検出モデルのみでは正確なグループ計測が困難となる要因が多数存在する。
本論文では、単純なベースライン手法（アプローチ01）から出発し、事前指定ROIによる空間制限、顔・骨格（ポーズ）検出の導入、エッジデバイス環境を考慮したWeb WorkerやONNX Runtime Web（YOLOv12 Nano）の活用、さらには複数フレーム情報を用いた人物トラッキング・通行人除外処理に至るまで、全13段階のアプローチによる逐次的な拡張と性能検証を行った。実験および各種アプローチの比較評価を通して、計算資源の制約下においてもOOM（Out of Memory）を防ぎつつ、待機列特有の幾何学的・時間的特徴を活用することで、高精度なグループ判定と堅牢な追跡処理が実現可能であることを示す。